%!TEX root = newmain.tex

To complete the translation, we encode the behaviour of a {\Stipula}
contract.  We use three key mechanisms:
%
\begin{enumerate}
\item Event declarations are collected into a \emph{dispatch table},
  recording enabled events and their execution times.
\item \emph{Symbolic execution} implicitly enumerates admissible
  scenarios via non-deter\-ministic function and event selection;
\item Time progression in a state {\Q} is bounded by a constant
  \stipula{max_tv}, defined as the maximum constant (plus one)
  occurring in time expressions, and tracked through the symbolic
  variable ${\tt u{\unds}Q} \in [0 .. \stipula{max_tv}]$. Advancing
  time beyond this bound is unnecessary, as no pending event can be
  executed past that deadline unless generated by a subsequent
  function invocation.
\end{enumerate}

  The dispatch table, introduced in Listing~\ref{tab:translate}, is an
  array \jml{DT} of size \jml{N}, where \jml{N} is the number of
  events in the contract.  The $i$-th entry corresponds to the $i$-th
  event
  $\stipula{now + t}_i\,\event\,\Qwithat_i\, \{ S_i \} \tostate\,
  \Qwithat_i'$.  The value \jml{DT[i]} is \jml{-1} if the event has
  not been created; otherwise, it stores the scheduled execution time
  $\jml{now + t}_i$, where \jml{now} denotes the time at which the
  event was generated.

  Events declared in a function are inserted in
  Listing~\ref{tab:translatefun} into the dispatch table by
  $\updateDT{\Q}{W}$ . If a function $\f$ declares events indexed by
  $\ev_{\jmath_1}, \ldots, \ev_{\jmath_{N_\f}}$ with time expressions
  $\jml{now + t}_{\jmath_i}$, then $\updateDT{\Q}{W}$ assigns the
  corresponding computed times to the entries
  \jml{DT[$\ev_{\jmath_i}$]}.  In other words, it updates exactly
  those positions associated with the events generated by
  $\f$. Entries of \jml{DT} not corresponding to events in $W$ remain
  unchanged, while the relevant entries are updated to
  $\jml{now + t}_{\jmath_i}$ using the current value of \jml{now}.
  The formalisation of this behaviour in {\JML} is done by the
  predicate \ensureDT$_\f$. Thus, the specification guarantees that
  precisely the events declared in $\f$ are scheduled, relative to the
  current time \jml{now}, and that no other dispatch table entries are
  affected.  (The formal definition of $\updateDT{\Q}{W}$ and
  {\ensureDT}$_\f$ is omitted, because it is straightforward.)


To obtain the possible behaviour of a \Stipula\ contract, the dispatch
table will be queried to identify one of the events that can be
enabled next. To this end, we introduce the auxiliary \Java\ method
%
\begin{center}
  \jml{int minTimeEntry(\{ev$_1$, $\ldots$ , ev$_m$\}) }\enspace.
\end{center}

% where {\tt ev}$_1$, $\cdots$ , {\tt ev}$_m$ are natural numbers
% encoding events of the contract.
The method returns the index in the dispatch table entry corresponding
to an event in \{{\tt ev}$_1$, $\ldots$ , {\tt ev}$_m$\} whose
associated time value is minimal, provided that such an event exists;
otherwise, it returns {\tt -1}. The method is non-deterministic, as
multiple events may share the same minimal time value in the dispatch
table. In this case, any one of these events may be selected, thereby
allowing the symbolic execution to explore all admissible scheduling
choices.  The implementation of the method \jml{minTimeEntry()} is
omitted. Instead, we provide its {\JML} specification that defines the
essential semantic properties required for verification, which is
sufficient for {\KeY} to reason symbolically about the method's
behaviour:

{\small
\begin{lstlisting}[language={[JML]Java},frame=none,numbers=none,mathescape=true,escapechar=']
/*@ public normal_behaviour // specification of int minTimeEntry({ev$_1$,$\ldots$, ev$_m$})
  @ ensures ( \exists i; 0 <= i < N; (i $\in$ {ev$_1$, $\ldots$, ev$_m$} && DT[i] >= now) )
            $\Rightarrow$ \result $\in$ {ev$_1$, $\ldots$, ev$_m$} && DT[\result] >= now && 
               \forall j; 0 <= j < N; (j $\in$ {ev$_1$, $\ldots$, ev$_m$} && DT[j] >= now) 
                                            $\Rightarrow$  DT[j] >= DT[\result] ;
  @ ensures ( \forall i; 0 <= i < N; i $\in$ {ev$_1$, $\ldots$, ev$_m$} $\Rightarrow$ DT[i] < now)
            $\Rightarrow$ \result == -1;
  @ assignable \nothing; 
  @*/
\end{lstlisting}
}

The first \jml{ensures} clause says that when there is an index in
\{\ev$_1$, $\ldots$, \ev$_m$\} such that its time value in \jml{DT} is
greater or equal to \jml{now}, then such an index is returned. There
can be more than one such index. This causes symbolic execution to
branch over all indices \jml{$\mathtt{\backslash}$result} such that
\jml{DT[$\mathtt{\backslash}$result] >= now}. The second \jml{ensures}
clause covers the case when no entry is found. In this case {\tt
  {\textbackslash}result} is \jml{-1}.

{\small
\begin{lstlisting}[frame=none,numbers=none,language={[JML]Java},mathescape=true,caption={The function $\genscenario{\C}{\Q}$; 
the ${\f}_1, \ldots, {\f}_n$ and ${\tt event\_1}, \ldots, {\tt event\_m}$ are the functions and events that start in {\Q} and end in $\Q_1, \ldots, \Q_n$ and 
$\Q_1', \ldots , \Q_m'$, respectively.},label={tab:fungenerate},float]
$\genscenario{\C}{\Q} \; \eqdef$ 
    int entry = minTimeEntry({ev$_1$,$\ldots$, ev$_m$}) ; 
    if ((entry == -1) || (DT[entry] > now && w_Q > 0)) { // call a function 
       if (entry == -1) now = now + u_Q ; else now = min(DT[entry]-1, now+u_Q) ;
       if (w_Q == 1) { $\f_1(\vect{\tt y_1}, \vect{\tt k_1})$ ; $\genscenario{\C}{\Q_1}$ }
       else if (w_Q == 2) { $\f_2(\vect{\tt y_2}, \vect{\tt k_2})$ ; $\genscenario{\C}{\Q_2}$ }
       ...
       else { $\f_n(\vect{\tt y_n}, \vect{\tt k_n})$ ; $\genscenario{\C}{\Q_n}$ } // max_Q == n, see Listing 1
    } else { // trigger an event
       now = DT[entry] ;
       if (entry == 1) { event_1() ; DT[entry] = -1 ; $\genscenario{\C}{\Q_1'}$ }
       else if (entry == 2) { event_2() ; DT[entry] = -1 ; $\genscenario{\C}{\Q_2'}$ }
       ...
       else { event_m() ; DT[entry] = -1 ; $\genscenario{\C}{\Q_m'}$ } // entry == m
    } 
\end{lstlisting}
}

  We can now complete the definition of \textsc{translate}
  (Listing~\ref{tab:translatefun}) by specifying the function
  $\genscenario{\C}{\Q}$ (Listing~\ref{tab:fungenerate}). Assume that
  in state $\Q$ functions ${\f}_1, \ldots, {\f}_n$ can be called and
  events ${\tt event}_1, \ldots, {\tt event}_m$ can be triggered. The
  {\Stipula} semantics is non-deterministic. To model this, we use two
  symbolic variables that were declared in
  Listing~\ref{tab:translate}: \jml{w_Q} $\in [\texttt{0}..\texttt{max\_Q}]$, for selecting a
  function in the state {\Q}, and \jml{u_Q} $\in [0..{\tt max\_tv}]$, for controlling time
  progression.  The behaviour is as follows:

  \begin{enumerate}
  \item \emph{Enabled events at current time}.  If
    \jml{minTimeEntry(...)} returns an index whose time equals
    \jml{now}, the corresponding event \emph{must} execute (events
    preempt functions). After execution, the event is removed from
    \jml{DT}. (This removal is actually unnecessary since the
    automaton is acyclic, it cannot occur again.)
  \item \emph{Pending future events.}  If one of the earliest
    scheduled events occurs at a time \jml{entry > now},
    then two cases arise: either time advances, but not until \jml{entry},
    i.e.\ at most to \jml{min(DT[entry]-1, now+u_Q)} and a function is
    executed (\jml{w_Q > 0}); otherwise, time advances \emph{exactly}
    to \jml{t} and the event is executed (\jml{w_Q = 0}).
  \item \emph{No pending events}.  If no event from $\Q$ remains to be
    scheduled, a function must eventually execute.  Time progression
    is bound by \jml{max_tv}, defined as the largest constant (plus
    one) appearing in time expressions. Assigning \jml{u_Q $\,\in\,$
      [0..max_tv]} ensures a finite, yet sound exploration of all
    relevant execution scenarios.
\end{enumerate}

Now that the translation is complete, it becomes possible to define and
verify meaningful properties for {\Stipula} contracts translated to
{\JML}/{\Java} by filling in the slots for \jml{OTHER PRE-CONDITION} and
\jml{THE POST-CONDITION}, respectively, in
Listing~\ref{tab:translate}. For example, one can prove with {\KeY}
the following property of the {\tt License} contract:

{\footnotesize
\begin{lstlisting}[language={[JML]Java},frame=none,numbers=none,mathescape=true,escapechar=']
requires Licensor_balance >= 0 && Licensee_balance > 0 &&
         Licensor_token > 0    && Licensee_token == 0 ;
ensures  (Licensor_balance == \old(Licensor_balance)+cost && Licensor_token == 0 
          && Licensee_balance == \old(Licensee_balance)-cost 
          && Licensee_token == \old(Licensor_token)
         ) || ( 
          Licensor_balance == \old(Licensor_balance) &&
          Licensor_token == \old(Licensor_token) && 
          Licensee_balance == \old(Licensee_balance) &&
          Licensee_token == \old(Licensee_token)
         ) ;
\end{lstlisting}
}

The contract expresses that, upon termination, either the licensing
transaction succeeded (the \stipula{Licensee} purchased the license
and the \stipula{Licensor} obtained the money) or the
\stipula{Licensee} is withdrawn and no \stipula{token} was delivered.

\begin{theorem}[Soundness]
  \label{thm.soundnessacyclic}
  Let $\C$ be an acyclic {\Stipula} contract
  % Let $\contract {\comma} {\tt 0}$ be the initial configuration of a
  % {\Stipula} acyclic contract {\C}
  and let
  \[
    \contract {\comma} {\tt 0} \lred{(\vect A,\, \vect{A_i}{:} \vect{v_i}{\,}^{i \in 1..\ell})}
    \lred{}^* \contract_1 {\comma} \Time_1 \lred{\mu_1} \lred{}^* 
    \contract_2 {\comma} \Time_2 \lred{\mu_2} \lred{}^* \cdots 
    % \contract_{n} {\comma} \Time_n
    \lred{\mu_{n+1}} \lred{}^* \contract_{n+1} {\comma} \Time_{n+1}
  \]
  where $\Q_i$ are the states of $\contract_i$, with $i \in 1..n+1$,
  $\Q_{n+1}$ is the final state of the underlying automaton, and
  $\mu_i$ are either function invocations
  $A_i{:}\f_i(\vect{v_i})[\vect{v_i'}]$ or events $\ev_i$. Then the
  above computation may be traced in the code of
  $\genscenario{\C}{\cdot}$ as follows. Consider
  $\contract_i {\comma} \Time_i \lred{\mu_i} \lred{}^* \contract_{i+1}
  {\comma} \Time_{i+1}$; if, in the code of $\genscenario{\C}{\Q_i}$,
  \begin{enumerate}
  \item \stipula{now} = $\Time_i$ if $\mu_i$ is an event;
    \stipula{now} $\leq \Time_i$ if $\mu_i$ is a function;
  \item if the multiset of pending events in $\contract_i$ is
    \[
      {\Time}_{\jmath_1} \,\event_{\ev_{\jmath_1}}\,
      \Qwithat_{\jmath_1}\, \{ \, S_{\jmath_1} \,\} \,\tostate\,
      \Qwithat_{\jmath_1}' ~|~ \cdots ~|~ {\Time}_{\jmath_\kappa}
      \event_{\ev_{\jmath_k}}\, \Qwithat_{\jmath_k}\, \{ \,
      S_{\jmath_k} \,\} \,\tostate\, \Qwithat_{\jmath_k}'
    \]
    and the dispatch table
    ${\tt DT}[\ev_{\jmath_i}] = {\Time}_{\jmath_i}$, with
    $i \in 1.. k$.
\end{enumerate}
Then, whenever $\mu_i = A_i{:}\f_i(\vect{v_i})[\vect{v_i'}]$ it is
possible to execute the method $\f_i(\vect{v_i},\vect{v_i'})$ and if
$\mu_i = \ev_{\jmath_i}$, it is possible to execute method
${\tt event}\_{\jmath_i}(~)$ and obtain a state where (1.) and (2.)
hold for $\contract_{i+1} \comma \Time_{i+1}$.
\end{theorem}

Therefore, every computation of an acyclic {\Stipula} contract is
faithfully represented within the {\Java} code generated by our
translation
%
% for each execution permitted by the {\Stipula} semantics of the
% contract, there exists a corresponding execution path in the
% translated {\Java}/{\JML} program.
(the proof is omitted, it is technical, but not difficult).
Consequently, when the program is analysed with {\KeY}, each legally
admissible behaviour of the original contract is systematically
explored. This guarantees that the verification process is sound with
respect to the contract's operational semantics in the acyclic
setting, as no admissible execution is omitted.  We further conjecture
that the translation is complete in the acyclic case, in the sense
that every behaviour of the generated {\Java} program corresponds to a
computation that is admissible in the original {\Stipula}
contract. Establishing such a result formally, however, would require
a formalisation of {\Java} semantics, which lies beyond the
scope of the present paper. 

%%% Local Variables:
%%% mode: latex
%%% TeX-master: "newmain"
%%% End:
